% =====================================================================
%  O GRUPO DE ROTACOES SO(3) E SUA ALGEBRA DE LIE
%  Nivel: inicio de mestrado em Fisica Teorica
%  Fonte: Computer Modern (padrao do pdfLaTeX, via codificacao T1/EC)
%
%  --------  COMO ALTERAR A FORMATACAO  --------------------------------
%  Toda a configuracao visual esta CONCENTRADA neste preambulo e
%  COMENTADA linha a linha. Para mudar margens, espacamento, recuo ou
%  tamanho de fonte, edite apenas os blocos indicados abaixo.
%
%  --------  MARCACOES EM VERMELHO  ------------------------------------
%  Tudo que foi ADICIONADO ou CORRIGIDO em relacao as notas originais
%  aparece em vermelho, iniciado pelo adereco  >  (ver comando \nota).
%  Para REMOVER todas as marcas de uma vez (deixar o texto uniforme),
%  troque a definicao de \nota por:   \renewcommand{\nota}[1]{#1}
% =====================================================================

\documentclass[12pt,a4paper]{article}   % 12pt: corpo do texto (ABNT)

% ---- Codificacao e idioma -------------------------------------------
\usepackage[T1]{fontenc}     % ativa a Computer Modern em codificacao T1
\usepackage[utf8]{inputenc}  % permite acentos digitados diretamente
\usepackage[brazil]{babel}   % hifenizacao e termos em portugues (ABNT)

% ---- Matematica ------------------------------------------------------
\usepackage{amsmath,amssymb}
\usepackage{bm}              % negrito matematico para vetores

% ---- Layout ABNT -----------------------------------------------------
% Margens ABNT: esquerda/superior 3 cm, direita/inferior 2 cm.
\usepackage[left=3cm,top=3cm,right=2cm,bottom=2cm]{geometry}

% Espacamento entre linhas 1,5 (ABNT para o corpo do texto).
\usepackage{setspace}
\onehalfspacing

% Recuo da primeira linha do paragrafo = 1,25 cm (ABNT).
\setlength{\parindent}{1.25cm}

% Recua tambem o primeiro paragrafo apos cada titulo de secao.
\usepackage{indentfirst}

% ---- Cores e titulos -------------------------------------------------
\usepackage{xcolor}
\usepackage[small,bf]{titlesec}  % titulos discretos; remova p/ o padrao

% =====================================================================
%  COMANDO DE MARCACAO EM VERMELHO
%  \nota{...} -> coloca o adereco vermelho ">" e escreve o trecho em
%  vermelho ate o fim. Funciona em texto e dentro de equacoes.
% =====================================================================
\definecolor{vermelho}{rgb}{0.80,0.00,0.00}
\newcommand{\nota}[1]{{\color{vermelho}%
  \ifmmode\blacktriangleright\,#1\else$\blacktriangleright$\,#1\fi}}

% ---- Abreviacoes uteis ----------------------------------------------
\newcommand{\R}{\mathbb{R}}
\newcommand{\so}{\mathfrak{so}}
\newcommand{\nvec}{\hat{\bm{n}}}          % versor do eixo
\newcommand{\Jvec}{\bm{J}}                % vetor de geradores
\newcommand{\ip}[2]{\langle #1,#2\rangle} % produto interno

% =====================================================================
\title{\textbf{O grupo de rota\c{c}\~oes $SO(3)$}\\[4pt]
       \large Eixo de rota\c{c}\~ao, geradores infinitesimais e o mapa exponencial}
\author{}          % preencha o autor, se desejar
\date{}            % preencha a data, se desejar

\begin{document}
\maketitle
\thispagestyle{empty}

\begin{center}
\fbox{\parbox{0.92\linewidth}{\small
\nota{Conven\c{c}\~ao: todo trecho em \textbf{vermelho}, iniciado pelo
adereco ``$\blacktriangleright$'', foi adicionado ou corrigido na
transcri\c{c}\~ao das notas manuscritas (passagens intermedi\'arias
explicitadas, um sinal corrigido e as rela\c{c}\~oes de comuta\c{c}\~ao
inclu\'idas). Para ocultar essas marcas, veja a instru\c{c}\~ao no topo
do arquivo-fonte.}}}
\end{center}

% =====================================================================
\section{O grupo $SO(n)$ e o eixo de rota\c{c}\~ao}
% =====================================================================

Definimos o \emph{grupo especial ortogonal} como o conjunto das
matrizes reais que preservam o produto interno e t\^em determinante
unit\'ario:
\begin{equation}
SO(n) := \bigl\{\, R\in M_n(\R) \;;\; \ip{Rx}{Ry}=\ip{x}{y}\ \
\forall\,x,y\in\R^{n},\ \ \det R = 1 \,\bigr\}.
\label{eq:defSO}
\end{equation}
A condi\c{c}\~ao $\ip{Rx}{Ry}=\ip{x}{y}$ equivale a $R^{\mathsf T}R=I$.
Se $R,S\in SO(n)$ ent\~ao $RS\in SO(n)$, e $SO(n)$ \'e um \emph{grupo de
Lie}.

\medskip
Mostremos que, para $n$ \'impar, toda rota\c{c}\~ao possui um
\textbf{eixo fixo}. Como $RR^{\mathsf T}=I$,
\begin{equation}
\det(R-I)=\det\!\bigl(R(I-R^{\mathsf T})\bigr)
        =\det R\;\det(I-R^{\mathsf T}).
\end{equation}
\nota{Usando $I-R^{\mathsf T}=-(R-I)^{\mathsf T}$ e
$\det A^{\mathsf T}=\det A$,}
\begin{equation}
\det(I-R^{\mathsf T})=(-1)^{n}\det(R-I).
\end{equation}
Com $\det R = 1$ e $n=2k+1$ \'impar, $(-1)^{n}=-1$, de modo que
\begin{equation}
\det(R-I) = -\det(R-I)\quad\Longrightarrow\quad \det(R-I)=0.
\end{equation}
Logo existe $v\neq 0$ com $Rv=v$; equivalentemente,
\begin{equation}
1\in\sigma(R)\qquad\forall\,R\in SO(n),\ n\ \text{\'impar}.
\end{equation}
\nota{Esse autovetor de autovalor $1$ \'e justamente a dire\c{c}\~ao
que a rota\c{c}\~ao deixa invariante --- o eixo.}

% =====================================================================
\section{Subespa\c{c}os invariantes}
% =====================================================================

Seja $W_R:=\ker(R-I)$ o subespa\c{c}o dos vetores fixos. Afirmamos que o
seu complemento ortogonal $W_R^{\perp}$ tamb\'em \'e invariante por $R$.
De fato, tomando $x\in W_R$ (logo $Rx=x$, e portanto $R^{\mathsf T}x=x$)
e $y\in W_R^{\perp}$,
\begin{equation}
\ip{x}{Ry}=\ip{R^{\mathsf T}x}{y}=\ip{x}{y}=0,
\end{equation}
o que mostra $Ry\in W_R^{\perp}$. Consequentemente o espa\c{c}o se
decomp\~oe como soma direta ortogonal,
\begin{equation}
V = W_R \oplus W_R^{\perp},
\label{eq:decomp}
\end{equation}
com $W_R$ n\~ao trivial quando $n$ \'e \'impar (supondo $R\neq I$).

% =====================================================================
\section{A dimens\~ao do eixo em $SO(3)$}
% =====================================================================

Especializando \eqref{eq:decomp} para $n=3$ e $R\neq I$, analisamos
$\dim W_R$:
\begin{itemize}\setlength{\itemsep}{2pt}
\item Se $\dim W_R = 3$, ent\~ao $W_R=V$, todo vetor \'e fixo e $R=I$
      --- exclu\'ido por hip\'otese.
\item Se $\dim W_R = 2$, ent\~ao $\dim W_R^{\perp}=1$; tome
      $0\neq y\in W_R^{\perp}$ com $Ry=\lambda y$. A preserva\c{c}\~ao
      da norma exige
      \begin{equation}
      \ip{Ry}{Ry}=\lambda^{2}\lVert y\rVert^{2}=\lVert y\rVert^{2}
      \;\Longrightarrow\; \lambda^{2}-1=0,\quad \lambda=\pm1.
      \end{equation}
      O valor $\lambda=+1$ recairia em $\dim W_R=3$; resta $\lambda=-1$,
      mas ent\~ao $\det R = (+1)(+1)(-1) = -1$, \nota{contrariando
      $\det R=1$. Logo $\dim W_R=2$ \'e imposs\'ivel.}
\end{itemize}
Portanto $\dim W_R = 1$: \nota{em $SO(3)$ toda rota\c{c}\~ao (n\~ao
trivial) fixa \emph{exatamente} uma reta --- o eixo de rota\c{c}\~ao.}

% =====================================================================
\section{Rota\c{c}\~oes elementares e forma em blocos}
% =====================================================================

Em duas dimens\~oes, $SO(2)$ \'e o grupo das rota\c{c}\~oes do plano,
\begin{equation}
SO(2)=\left\{
\begin{pmatrix} \cos\theta & -\sin\theta \\ \sin\theta & \cos\theta
\end{pmatrix};\ \theta\in[0,2\pi)\right\}.
\end{equation}
Escolhendo uma base ortonormal cujo primeiro vetor aponte ao longo do
eixo fixo, uma rota\c{c}\~ao de $SO(3)$ assume a forma em blocos
\begin{equation}
R=\begin{pmatrix} 1 & \bm{0}^{\mathsf T}\\[2pt] \bm{0} & \tilde R
\end{pmatrix},\qquad \tilde R\in SO(2),
\end{equation}
que age trivialmente sobre o eixo e como rota\c{c}\~ao plana no
complemento. As tr\^es rota\c{c}\~oes elementares em torno dos eixos
coordenados s\~ao
\begin{align}
R_1(\varphi)&=\begin{pmatrix}
1 & 0 & 0\\ 0 & \cos\varphi & -\sin\varphi\\ 0 & \sin\varphi & \cos\varphi
\end{pmatrix},
&
R_2(\varphi)&=\begin{pmatrix}
\cos\varphi & 0 & \nota{+\sin\varphi}\\ 0 & 1 & 0\\ -\sin\varphi & 0 & \cos\varphi
\end{pmatrix},
\\[4pt]
R_3(\varphi)&=\begin{pmatrix}
\cos\varphi & -\sin\varphi & 0\\ \sin\varphi & \cos\varphi & 0\\ 0 & 0 & 1
\end{pmatrix}.
\end{align}
\nota{Corre\c{c}\~ao: nas notas, o elemento superior direito de
$R_2$ aparecia como $-\sin\varphi$; o sinal correto \'e $+\sin\varphi$,
como se v\^e ao derivar $R_2$ e comparar com o gerador $J_2$ obtido
adiante.}

% =====================================================================
\section{Geradores infinitesimais e a \'algebra $\so(3)$}
% =====================================================================

Derivando cada rota\c{c}\~ao elementar na origem obtemos os
\emph{geradores infinitesimais} $J_a := \dfrac{dR_a}{d\varphi}\Big|_{\varphi=0}$:
\begin{equation}
J_1=\begin{pmatrix}0&0&0\\0&0&-1\\0&1&0\end{pmatrix},\quad
J_2=\begin{pmatrix}0&0&1\\0&0&0\\-1&0&0\end{pmatrix},\quad
J_3=\begin{pmatrix}0&-1&0\\1&0&0\\0&0&0\end{pmatrix}.
\end{equation}
Em componentes, $(J_a)_{ij}=-\varepsilon_{aij}$; s\~ao uma base do
espa\c{c}o das matrizes antissim\'etricas $3\times3$. Assim
\begin{equation}
\so(3)=\mathrm{span}\{J_1,J_2,J_3\},
\end{equation}
que, munido do comutador $[\,\cdot\,,\cdot\,]$, \'e a \emph{\'algebra de
Lie} do grupo. \nota{Sua estrutura \'e fixada pelas rela\c{c}\~oes de
comuta\c{c}\~ao}
\begin{equation}
\nota{[J_a,J_b]=\varepsilon_{abc}\,J_c,}
\label{eq:comm}
\end{equation}
\nota{que reproduzem, a menos de fatores, a \'algebra do momento angular.}
Cada rota\c{c}\~ao elementar \'e recuperada pela exponencial,
$R_a(\varphi)=\exp(\varphi J_a)$.

% =====================================================================
\section{O mapa exponencial e a pertin\^encia a $SO(3)$}
% =====================================================================

\noindent\textbf{Proposi\c{c}\~ao.} \emph{Seja $\theta\in[0,\pi]$ e
$\nvec\in\R^{3}$ com $\lVert\nvec\rVert=1$. A rota\c{c}\~ao de \^angulo
$\theta$ em torno de $\nvec$ \'e}
\begin{equation}
R(\theta,\nvec)=\exp\!\bigl(\theta\,\nvec\cdot\Jvec\bigr),
\qquad
\nvec\cdot\Jvec=\sum_i n_i J_i
=\begin{pmatrix}0&-n_3&n_2\\ n_3&0&-n_1\\ -n_2&n_1&0\end{pmatrix}.
\label{eq:nJ}
\end{equation}
\nota{Note que $(\nvec\cdot\Jvec)\,u=\nvec\times u$: a matriz
\eqref{eq:nJ} \'e a do produto vetorial por $\nvec$.}

\medskip
A matriz $\nvec\cdot\Jvec$ \'e antissim\'etrica,
$(\nvec\cdot\Jvec)^{\mathsf T}=-\,\nvec\cdot\Jvec$. Logo
\begin{equation}
\bigl[\exp(\theta\,\nvec\cdot\Jvec)\bigr]^{\mathsf T}
=\exp\!\bigl(\theta\,(\nvec\cdot\Jvec)^{\mathsf T}\bigr)
=\exp(-\theta\,\nvec\cdot\Jvec)
=\bigl[\exp(\theta\,\nvec\cdot\Jvec)\bigr]^{-1},
\end{equation}
isto \'e, $\exp(\theta\,\nvec\cdot\Jvec)\in O(3)$. \nota{Al\'em disso,
como o tra\c{c}o de uma matriz antissim\'etrica \'e nulo,}
\begin{equation}
\det\exp(\theta\,\nvec\cdot\Jvec)
=\exp\!\bigl(\theta\,\mathrm{Tr}(\nvec\cdot\Jvec)\bigr)=e^{0}=1,
\end{equation}
e portanto $R(\theta,\nvec)\in SO(3)$.

% =====================================================================
\section{F\'ormula de Rodrigues e parametriza\c{c}\~ao de $SO(3)$}
% =====================================================================

A a\c{c}\~ao de $R(\theta,\nvec)$ sobre um vetor arbitr\'ario \'e dada
pela \emph{f\'ormula de Rodrigues}
\begin{equation}
e^{\theta\,\nvec\cdot\Jvec}\,u
= u + (1-\cos\theta)\,\nvec\times(\nvec\times u) + \sin\theta\,(\nvec\times u).
\label{eq:rodrigues}
\end{equation}
Aplicada ao pr\'oprio eixo, com $\nvec\times\nvec=\bm 0$,
\begin{equation}
e^{\theta\,\nvec\cdot\Jvec}\,\nvec = \nvec,
\end{equation}
confirmando que $\nvec$ permanece fixo. Para $u\perp\nvec$, usamos
\nota{a identidade $\nvec\times(\nvec\times u)=\nvec(\nvec\cdot u)-(\nvec\cdot\nvec)\,u$} em \eqref{eq:rodrigues}:
\begin{align}
e^{\theta\,\nvec\cdot\Jvec}\,u
&= u + (1-\cos\theta)\bigl[\nvec(\nvec\cdot u)-u\bigr]+\sin\theta\,(\nvec\times u)
\nonumber\\
&= u\cos\theta + \sin\theta\,(\nvec\times u),
\end{align}
pois $\nvec\cdot u=0$ e $\lVert\nvec\rVert=1$. Analogamente,
\begin{equation}
e^{\theta\,\nvec\cdot\Jvec}\,(\nvec\times u)
= -\sin\theta\,u + \cos\theta\,(\nvec\times u).
\end{equation}
\nota{Ou seja, no plano ortogonal a $\nvec$ o operador age exatamente
como uma rota\c{c}\~ao de \^angulo $\theta$}, o que justifica a
interpreta\c{c}\~ao geom\'etrica. Conclu\'imos que
\begin{equation}
\boxed{\;SO(3)=\bigl\{\exp(\theta\,\nvec\cdot\Jvec)\;;\;
\theta\in[0,\pi],\ \nvec\in\R^{3},\ \lVert\nvec\rVert=1\bigr\}.\;}
\end{equation}
\nota{A exponencial \'e sobrejetiva sobre $SO(3)$ porque o grupo \'e
conexo e compacto; o intervalo $[0,\pi]$ (e n\~ao $[0,2\pi]$) basta
porque $R(\theta,\nvec)=R(2\pi-\theta,-\nvec)$.}

% =====================================================================
\section*{Refer\^encias}   % ABNT NBR 6023 (secao nao numerada)
% =====================================================================
\begin{flushleft}
\setlength{\parindent}{0pt}

\nota{HALL, B. C. \textbf{Lie Groups, Lie Algebras, and
Representations}. 2. ed. New York: Springer, 2015.}

\vspace{4pt}
\nota{GILMORE, R. \textbf{Lie Groups, Physics, and Geometry}.
Cambridge: Cambridge University Press, 2008.}

\vspace{4pt}
\nota{TUNG, W.-K. \textbf{Group Theory in Physics}. Singapore: World
Scientific, 1985.}
\end{flushleft}

\end{document}
